\documentclass[aspectratio=169, 11pt]{beamer}

\usepackage[utf8]{inputenc}
\usepackage[T1]{fontenc}
\usepackage[ngerman]{babel}
\usepackage{helvet}
\renewcommand{\familydefault}{\sfdefault}
\usepackage{listings}
\usepackage{xcolor}

% -------------------------------------------------------
% Theme: Metropolis (modern, serifenlos, minimal)
% -------------------------------------------------------
\usetheme[
  progressbar=frametitle,
  block=fill,
  numbering=fraction,
]{metropolis}

% Farben -- gedämpfte Palette, kein grelles Blau
\definecolor{darkslate}{HTML}{0F172A}   % fast schwarz, für Titel + Text
\definecolor{slategray}{HTML}{475569}   % mittelgrau für Akzente
\definecolor{lightbg}{HTML}{F8FAFC}     % sehr helles Grau für Hintergrund
\definecolor{rulecolor}{HTML}{CBD5E1}   % Trennlinien
\definecolor{hinweisbg}{HTML}{F1F5F9}
\definecolor{mutedtext}{HTML}{64748B}
\definecolor{codeblü}{HTML}{3B82F6}
\definecolor{codegreen}{HTML}{16A34A}
\definecolor{codecomment}{HTML}{94A3B8}

\setbeamercolor{normal text}{fg=darkslate, bg=white}
\setbeamercolor{alerted text}{fg=slategray}
\setbeamercolor{frametitle}{fg=white, bg=darkslate}
\setbeamercolor{progress bar}{fg=slategray, bg=rulecolor}
\setbeamercolor{block title}{fg=white, bg=darkslate}
\setbeamercolor{block body}{fg=darkslate, bg=lightbg}
\setbeamercolor{title separator}{fg=slategray}
\setbeamercolor{footline}{fg=mutedtext}

% Titelfolie
\setbeamercolor{title}{fg=white}
\setbeamercolor{subtitle}{fg=white!70!darkslate}
\setbeamercolor{author}{fg=white!70!darkslate}
\setbeamercolor{institute}{fg=white!70!darkslate}
\setbeamercolor{date}{fg=white!60!darkslate}
\setbeamercolor{background canvas}{bg=white}

% Schriftgrössen
\setbeamerfont{title}{size=\large, series=\bfseries}
\setbeamerfont{subtitle}{size=\normalsize, series=\mdseries}
\setbeamerfont{frametitle}{size=\normalsize, series=\bfseries}
\setbeamerfont{author}{size=\small}
\setbeamerfont{date}{size=\small}
\setbeamerfont{block title}{size=\small, series=\bfseries}
\setbeamerfont{block body}{size=\small}

% Aufzählungszeichen: simpler Strich
\setbeamertemplate{itemize item}{\small\raise1pt\hbox{--}}
\setbeamertemplate{itemize subitem}{\scriptsize\raise1pt\hbox{--}}

% -------------------------------------------------------
% Code-Stil
% -------------------------------------------------------
\lstdefinestyle{pyslide}{
  language=Python,
  basicstyle=\ttfamily\scriptsize,
  keywordstyle=\color{codeblü}\bfseries,
  commentstyle=\color{codecomment}\itshape,
  stringstyle=\color{codegreen},
  backgroundcolor=\color{hinweisbg},
  frame=none, numbers=none,
  showstringspaces=false, breaklines=trü, tabsize=4,
  xleftmargin=6pt, xrightmargin=6pt,
  aboveskip=4pt, belowskip=4pt,
  literate={ä}{{\"a}}1{ö}{{\"o}}1{ü}{{\"u}}1{ß}{{\ss}}1,
}

% Hinweiskasten (vor Abgabe entfernen)
\newcommand{\hinweis}[1]{%
  \vspace{4pt}%
  {\scriptsize\color{mutedtext}\textit{#1}}%
}

% -------------------------------------------------------
\title{Bottleneck-Analyse im Karosseriebau}
\subtitle{Eure Fragestellung in einem Satz}
\author{Vorname Name \quad Vorname Name}
\institute{Datenanalyse im Produktionsumfeld}
\date{Sommersemester 2025}

% -------------------------------------------------------
\begin{document}

\maketitle

% ----------------------------------------------------------------
\begin{frame}{Fragestellung}

  \begin{block}{Unsere Frage}
    \textit{[Eine einzige, klar formulierte Frage.]}\\
    {\small Beispiel: Ist die erhöhte Taktzeit an Station S04 auf einen
    einzelnen Roboter zurückzuführen -- oder ist sie typabhängig?}
  \end{block}

  \vspace{8pt}
  \textbf{Warum relevant für die Produktion:}\\[4pt]
  \textit{[1--2 Sätze. Welche Entscheidung hängt von der Antwort ab?
  Konkret: Wartungsplanung, Kapazität, Reihenfolgeplanung?]}

  \vspace{8pt}
  \hinweis{Diese Folie in 60 Sekunden. Kein Kontext -- direkt zur Frage.}

\end{frame}

% ----------------------------------------------------------------
\begin{frame}{Daten \& Vorgehen}

  \begin{columns}[T, onlytextwidth]
    \begin{column}{0.47\textwidth}
      \textbf{Datenbasis}
      \begin{itemize}
        \item \textit{[Anzahl Karosserien, Zeitraum]}
        \item \textit{[Stationen, Roboter, Typen]}
        \item Echte Daten vertraulich --\\
          Analyse auf Beispieldatensatz entwickelt
      \end{itemize}
    \end{column}

    \begin{column}{0.47\textwidth}
      \textbf{Analytischer Weg}
      \begin{itemize}
        \item \textit{[Schritt 1: z.\,B. Verteilungsanalyse]}
        \item \textit{[Schritt 2: z.\,B. Gruppenvergleich]}
        \item \textit{[Schritt 3: z.\,B. lineares Modell]}
      \end{itemize}

      \vspace{6pt}
      \textbf{Warum so?}\\
      {\small\textit{[1 Satz: was hat euch zu diesem
      Weg geführt?]}}
    \end{column}
  \end{columns}

  \hinweis{Methoden nicht erklären -- zeigen, dass ihr wisst was ihr getan habt.}

\end{frame}

% ----------------------------------------------------------------
\begin{frame}{Kernergebnis}
  % Falls zwei aufeinander aufbaünde Befunde: diese Folie duplizieren.

  \begin{columns}[T, onlytextwidth]
    \begin{column}{0.58\textwidth}
      % \includegraphics[width=\linewidth]{plots/eür_plot.png}
      \begin{center}
        \textcolor{rulecolor}{\rule{\linewidth}{0.4pt}}
        \vspace{1.2cm}
        {\small\textcolor{mutedtext}{Grafik einfügen}}
        \vspace{1.2cm}
        \textcolor{rulecolor}{\rule{\linewidth}{0.4pt}}
      \end{center}
    \end{column}

    \begin{column}{0.38\textwidth}
      \textbf{Was wir sehen}\\[6pt]
      {\small\textit{[Kernaussage in 1--2 Sätzen.
      Konkret: welcher Wert, welcher Unterschied,
      welche Station oder welcher Roboter?]}}

      \vspace{10pt}
      \textbf{Wie sicher?}\\[6pt]
      {\small\textit{[N, Konfidenzintervall, oder:
      ,,Muster ist konsistent über alle Typen``]}}
    \end{column}
  \end{columns}

  \hinweis{Grafik nicht vorlesen -- erklären. Was sieht man, was bedeutet es?}

\end{frame}

% ----------------------------------------------------------------
\begin{frame}{Implikation für die Produktion}

  \vspace{4pt}
  \textbf{Was das bedeutet:}\\[8pt]
  {\normalsize\textit{[Konkreter Satz in Produktionssprache -- nicht
  ,,die Taktzeit ist höher``, sondern z.\,B.:
  ,,Solange Roboter R02 nicht gewartet wird, verliert diese Linie
  pro Schicht X Takteinheiten -- unabhängig davon, wie
  gut alle anderen Stationen laufen.``]}}

  \vspace{14pt}
  \textbf{Frage an den Konzern:}\\[6pt]
  {\small\textit{[Was müsste der Konzern als nächstes prüfen,
  entscheiden oder bereitstellen?
  Richtet diese Frage direkt an die Anwesenden.]}}

  \hinweis{Das ist der Kern der Zwischenpräsi.}

\end{frame}

% ----------------------------------------------------------------
\begin{frame}{Grenzen \& Ausblick}

  \begin{columns}[T, onlytextwidth]
    \begin{column}{0.47\textwidth}
      \textbf{Was diese Analyse nicht erfasst}
      \begin{itemize}
        \item \textit{[z.\,B. keine Kausalität,\\
          nur Korrelation]}
        \item \textit{[z.\,B. Schichteffekte nicht\\
          prüfbar, Zeitstempel fehlen]}
        \item \textit{[z.\,B. nur eine Station\\
          tief analysiert]}
      \end{itemize}
    \end{column}

    \begin{column}{0.47\textwidth}
      \textbf{Nächste sinnvolle Schritte}
      \begin{itemize}
        \item \textit{[Was würde die Analyse\\
          ermöglichen, die jetzt fehlt?]}
        \item \textit{[Was müsste der Konzern\\
          dafür bereitstellen?]}
      \end{itemize}
    \end{column}
  \end{columns}

  \vspace{10pt}
  \hinweis{Grenzen zu kennen ist keine Schwäche --
  es zeigt, dass ihr den Scope eurer Arbeit verstanden habt.}

\end{frame}

\end{document}
